Software systems tend to grow as they evolve, yet how this growth manifests in the base of a long-lived Linux distribution has not been characterised systematically. This thesis provides a reproducible, release-wise characterisation of the persistent core of Debian's minbase installation, the minimal package set that \texttt{debootstrap} resolves for each release: the 25 source packages present in this set in at least 11 of the 13 stable releases from Potato (2000) to Trixie (2025). A containerised pipeline downloads each release's package indices and source archives and computes five metrics per release: installed size, dependency structure, lines of code, per-function cyclomatic complexity, and distribution-specific patch volume. The results describe growth by accretion. The persistent core grows roughly 5.3-fold in lines of code and roughly fourfold in installed size, while the median per-function cyclomatic complexity remains at 3 in every release and the median per-package dependency count at 1 in every release after Potato; the composition of the minbase changes chiefly through packaging structure rather than capability. The patch-volume metric is comparable only from the 3.0 (quilt) source format onward and declines within that window. The study is exploratory-descriptive and asserts no causal or correlational relationship between the dimensions. The pipeline and an interactive companion application are made publicly available to enable reproduction and extension of the analysis.

\vspace{6em} 

\noindent\textbf{Keywords:} Debian, software evolution, software bloat, software metrics, cyclomatic complexity, mining software repositories
